\documentclass[12pt,a4paper,UTF8]{ctexart}
\usepackage[a4paper,top=22mm,bottom=22mm,left=22mm,right=22mm,headheight=15pt]{geometry}
\newcommand{\Needspace}[1]{\par\begingroup\dimen0=#1\relax\dimen2=\pagegoal\advance\dimen2 by-\pagetotal\ifdim\dimen0>\dimen2\newpage\fi\endgroup}
\newcommand{\needspace}[1]{\par\begingroup\skip0=\lastskip\vskip-\skip0\dimen0=#1\relax\advance\dimen0 by\skip0\vskip\dimen0\penalty-100\vskip-\dimen0\vskip\skip0\penalty9999\endgroup}
\usepackage{amsmath,amssymb,mathtools,bm}
\usepackage{graphicx,xcolor,booktabs,longtable,array,enumitem,fancyhdr,caption,float,microtype,setspace}
\usepackage[hidelinks,unicode]{hyperref}
\usepackage[most]{tcolorbox}
\definecolor{ContestBlue}{HTML}{17365D}
\hypersetup{pdftitle={第073套},pdfauthor={}}
\setstretch{1.24}
\setlength{\parindent}{2em}
\setlength{\parskip}{0.35em}
\setlength{\emergencystretch}{2em}
\setlist{itemsep=0.35em,topsep=0.4em}
\captionsetup{font=small,labelfont=bf,justification=centering}
\setcounter{secnumdepth}{0}
\pagestyle{fancy}
\fancyhf{}
\fancyhead[L]{\small 第073套}
\fancyfoot[C]{\thepage}
\renewcommand{\headrulewidth}{0.4pt}
\ctexset{
  section={format=\Large\bfseries\color{ContestBlue},beforeskip=1.2em,afterskip=0.7em},
  subsection={format=\large\bfseries\color{ContestBlue},beforeskip=1.0em,afterskip=0.5em},
  subsubsection={format=\normalsize\bfseries,beforeskip=0.8em,afterskip=0.35em}
}
\usepackage{calc}
\newcounter{none}
\let\originalsubsection\subsection
\renewcommand{\subsection}{\needspace{6\baselineskip}\originalsubsection}
\makeatletter
\newcommand{\tightlist}{\setlength{\itemsep}{0pt}\setlength{\parskip}{0pt}}
\newcommand{\maxwidth}{\ifdim\Gin@nat@width>\linewidth\linewidth\else\Gin@nat@width\fi}
\newcommand{\maxheight}{\ifdim\Gin@nat@height>0.72\textheight 0.72\textheight\else\Gin@nat@height\fi}
\newcommand{\pandocbounded}[1]{\begingroup\setkeys{Gin}{width=\maxwidth,height=\maxheight,keepaspectratio}#1\endgroup}
\makeatother
\begin{document}
\begin{titlepage}
  \centering
  \vspace*{58mm}
  {\Huge\bfseries\color{ContestBlue} 第073套\par}
  \vfill
  \vspace*{28mm}
\end{titlepage}
\section{第 1 题}\label{ux7b2c-1-ux9898}

两块大接地金属板构成二面角 \(\theta\)。一根长 \(L\)、均匀带电总量 \(Q\)
的细杆平行于二面角棱线，位于角平分面内，距顶点（棱线）为 \(r\)，且
\(L\gg r\)。求杆所受静电力，分别讨论 \(\theta=180^\circ\) 与
\(\theta=18^\circ\)。

\begin{figure}[H]
\centering
\includegraphics[width=0.76\linewidth,height=\textheight,keepaspectratio,alt={图1-1　接地金属二面角内沿棱线方向的长带电杆}]{../figures/question-01.pdf}
\caption*{图1-1　接地金属二面角内沿棱线方向的长带电杆}
\end{figure}

\newpage

\section{第 2 题}\label{ux7b2c-2-ux9898}

若干等长轻质刚杆铰接成图示可变形支架，左端固定，右端由支撑脚支起；顶点和左侧指定位置分别受竖直负载
\(G,P\)。图示中央下横杆的轴力记为 \(N\)。

\begin{enumerate}
\def\labelenumi{\arabic{enumi}.}
\tightlist
\item
  支撑脚与地面无摩擦，求 \(N\)。
\item
  支撑脚与地面摩擦因数 \(\mu\le1\)，求 \(N\) 的允许范围。
\end{enumerate}

\begin{figure}[H]
\centering
\includegraphics[width=0.695\linewidth,height=\textheight,keepaspectratio,alt={图2-1　铰接杆支架、两处负载与中央横杆}]{../figures/question-02.pdf}
\caption*{图2-1　铰接杆支架、两处负载与中央横杆}
\end{figure}

\newpage

\section{第 3 题}\label{ux7b2c-3-ux9898}

厚度为 \(l\) 的均匀长板放在两根轴距 \(2l\)
的高速反向滚筒上。两滚筒正上方各有固定柱，平衡时与板上表面虚接触。板与滚筒、固定柱间摩擦因数均为
\(\mu=1/2\)。初始板质心在两滚筒中点，以速度 \(v_0\) 向右。

\begin{enumerate}
\def\labelenumi{\arabic{enumi}.}
\tightlist
\item
  板质心到右滚筒正上方 \(x=l\) 时恰静止，求 \(v_0\)。
\item
  求板第一次回到中点的速度 \(u_1\)。
\item
  求充分久后板从左向右回到中点的速度 \(u_2\)。
\end{enumerate}

\begin{figure}[H]
\centering
\includegraphics[width=0.84\linewidth,height=\textheight,keepaspectratio,alt={图3-1　反向滚筒、上方固定柱与横向运动长板}]{../figures/question-03.pdf}
\caption*{图3-1　反向滚筒、上方固定柱与横向运动长板}
\end{figure}

\newpage

\section{第 4 题}\label{ux7b2c-4-ux9898}

两根平行光滑导轨间距 \(l\)，导轨面倾角
\(\theta=\pi/6\)，底部以导线连接。空间有竖直向下匀强磁场 \(B\)。质量
\(m\) 的水平导体棒与导轨构成零电阻回路，自感

\[
L(x)=kx,\qquad
k=\frac{B^2l^2}{4mg}.
\]

初态棒静止于 \(x_0\)，电流为零，且运动区间内
\(x\gg l\)。沿斜面向上给瞬时冲量 \(I_0\)。

\begin{enumerate}
\def\labelenumi{\arabic{enumi}.}
\tightlist
\item
  求瞬时速度。
\item
  要使棒此后始终贴在导轨上运动，求 \(I_0\) 范围。
\end{enumerate}

\begin{figure}[H]
\centering
\includegraphics[width=0.72\linewidth,height=\textheight,keepaspectratio,alt={图4-1　倾斜超导导轨、水平导体棒与竖直磁场}]{../figures/question-04.pdf}
\caption*{图4-1　倾斜超导导轨、水平导体棒与竖直磁场}
\end{figure}

\newpage

\section{第 5 题}\label{ux7b2c-5-ux9898}

光滑水平面上，一根长 \(2l\) 的轻杆两端固定质量均为 \(m\)
的小球，电量分别为 \(\pm q\)。匀强磁场 \(B\) 竖直向下。初始杆中心速度
\(v_0\) 垂直于杆，角速度 \(\omega_0\)。以初始杆中心为原点，初始杆从
\(-q\) 指向 \(+q\) 为 \(y\) 轴，\(v_0\) 方向为 \(x\) 轴。杆从 \(-q\)
指向 \(+q\) 的方向从 \(+y\) 轴逆时针转过的有向角为
\(\theta\)，\(\omega=\dot\theta\)，初态 \(v_0,\omega_0\ge0\)。

\begin{enumerate}
\def\labelenumi{\arabic{enumi}.}
\tightlist
\item
  求杆中心速度 \(v_x(\theta),v_y(\theta)\) 与 \(\omega(\theta)\)。
\item
  求杆扫过区域的最小带宽 \(\Delta h(\omega_0)\)。
\end{enumerate}

\begin{figure}[H]
\centering
\includegraphics[width=0.74\linewidth,height=\textheight,keepaspectratio,alt={图5-1　异号带电双球刚性杆在垂直磁场中的初态}]{../figures/question-05.pdf}
\caption*{图5-1　异号带电双球刚性杆在垂直磁场中的初态}
\end{figure}

\newpage

\section{第 6 题}\label{ux7b2c-6-ux9898}

土豆大炮的炮管截面积为 \(A\)。发射瞬间土豆静止，土豆与封闭端之间气体体积
\(V_0\)、压强
\(P_0>P_{\rm atm}\)。气体为二原子理想气体，快速绝热膨胀，无摩擦、无泄漏。炮管总长度
\(L\) 从封闭端量到管口，忽略土豆自身的纵向长度。改变 \(L\)：

\begin{enumerate}
\def\labelenumi{\arabic{enumi}.}
\tightlist
\item
  求土豆出管最大动能。
\item
  求实现最大动能的 \(L\)。
\end{enumerate}

\begin{figure}[H]
\centering
\includegraphics[width=0.82\linewidth,height=\textheight,keepaspectratio,alt={图6-1　绝热气体推动土豆并克服外界大气压}]{../figures/question-06.pdf}
\caption*{图6-1　绝热气体推动土豆并克服外界大气压}
\end{figure}

\newpage

\section{第 7 题}\label{ux7b2c-7-ux9898}

焦距 \(f\) 的凸透镜沿光轴切成上下两半，拉开后以厚度 \(\delta\)
的不透光板填充中缝。轴上点光源距透镜 \(p>f\)，波长 \(\lambda\)。屏距透镜
\(H\)。已知

\[
f=10\,\mathrm{cm},\quad p=20\,\mathrm{cm},\quad
\delta=1\,\mathrm{mm},\quad
\lambda=0.5\,\mu\mathrm m,\quad H=50\,\mathrm{cm}.
\]

求屏上可见干涉条纹数。

\begin{figure}[H]
\centering
\includegraphics[width=0.72\linewidth,height=\textheight,keepaspectratio,alt={图7-1　分割并拉开的凸透镜形成双像干涉}]{../figures/question-07.pdf}
\caption*{图7-1　分割并拉开的凸透镜形成双像干涉}
\end{figure}

\newpage

\section{第 8 题}\label{ux7b2c-8-ux9898}

静止离子可吸收波长 \(\lambda_0=580.000\,\mathrm{nm}\) 的光。离子以 \(0\)
到 \(v_2=5000\,\mathrm{m\,s^{-1}}\) 朝激光光源运动。

\begin{enumerate}
\def\labelenumi{\arabic{enumi}.}
\tightlist
\item
  求激发全部离子所需激光波长范围，并画出荧光光子数的波长分布。
\item
  离子先经电压 \(V\) 加速，求速度区间宽度与 \(V\)
  的关系，判断变宽或变窄。
\item
  荷质比 \(q/m=5.00\times10^6\,\mathrm{C\,kg^{-1}}\)，有两跃迁静止波长
  \(580.000\,\mathrm{nm}\)、\(580.001\,\mathrm{nm}\)。求无预加速时两扫描范围；再求使两范围不重叠的最小加速电压。
\end{enumerate}

\begin{figure}[H]
\centering
\includegraphics[width=0.88\linewidth,height=\textheight,keepaspectratio,alt={图8-1　逆向激光吸收测量离子束速度分布}]{../figures/question-08.pdf}
\caption*{图8-1　逆向激光吸收测量离子束速度分布}
\end{figure}
\end{document}
